\section{Controlled PLAID Comparison}
\label{sec:results-plaid}

PLAID is a purpose-built engine for ColBERT-style late interaction and is
therefore a more relevant system baseline than reducing MaxSim to ordinary
single-vector search. The earlier project comparison was invalid because PLAID
could rescore up to 8{,}192 documents while the IVF pipelines reranked only 50,
used a different timing boundary, and came from a different runtime stack. This
section replaces it with measurements made in the same Python process,
interleaving schedule, CPU affinity, thread budget, query split, and online
boundary as the IVF arms of \cref{sec:results-ivf}.

The controlled runner uses PyLate~1.5.0 and CPU
\texttt{fast-plaid}~1.3.0.290 with four-bit residual compression
($\texttt{nbits}=4$), four k-means iterations, and Triton disabled. The public
API returns final ranked documents but does not expose the realized candidate
pool. Consequently, \texttt{n\_full\_scores} is reported as a configured
full-scoring upper bound, not as an observed candidate count or a quantity
strictly equivalent to IVF's reranked $C$.

\subsection{Predeclared selection and held-out results}

The first 10 SciFact queries were used only for parameter selection. Across the
initial grid, $(\texttt{n\_ivf\_probe},\texttt{n\_full\_scores})=(8,100)$ was
the fastest point above the predeclared validation exact-recall@10 threshold of
0.85 (observed 0.89), while $(8,400)$ achieved the maximum observed validation
recall (0.92). Increasing the full-score budget to 800, 1{,}600, 3{,}200, or
5{,}183 did not improve that validation value. The two selected points were
then frozen before evaluation on queries 10--49 of SciFact and NFCorpus.

\begin{table}[tb]
  \centering
  \caption{Held-out results for the two validation-selected PLAID operating
    points (40 queries per dataset; five interleaved repetitions). Exact R@10
    is recovery against compiled exhaustive MaxSim. ``Budget ratio'' is
    \texttt{n\_full\_scores} divided by corpus size and is an upper-bound
    configuration ratio, not an observed rerank ratio.}
  \label{tab:plaid-selected}
  \small
  \setlength{\tabcolsep}{4.5pt}
  \begin{tabular}{@{}l S[table-format=2.4] S[table-format=2.4] S[table-format=1.4] S[table-format=1.4] S[table-format=1.4]@{}}
    \toprule
    Dataset / arm & {Median (\si{\second})} & {p95 (\si{\second})} & {Exact R@10} & {Qrels R@10} & {Budget ratio} \\
    \midrule
    SciFact, $F{=}100$  & 14.5497 & 14.7523 & 0.2875 & 0.3250 & 0.0193 \\
    SciFact, $F{=}400$  & 16.6456 & 16.8475 & 0.4375 & 0.3750 & 0.0772 \\
    NFCorpus, $F{=}100$ &  5.6825 &  6.0592 & 0.3100 & 0.1434 & 0.0275 \\
    NFCorpus, $F{=}400$ &  7.2505 &  7.3043 & 0.3950 & 0.1043 & 0.1101 \\
    \bottomrule
  \end{tabular}
\end{table}

The held-out values in \cref{tab:plaid-selected} are much lower than the
validation values: SciFact exact recall@10 falls from 0.89/0.92 to
0.2875/0.4375. The same frozen points reach only 0.3100/0.3950 on NFCorpus.
This is a real split shift rather than an embedding-normalization mismatch:
the source token vectors already have unit norm to numerical tolerance, and
the release index reproduces the validation scores before it is used on the
held-out split. The selected PLAID arms are also slower than the compiled exact
and IVF arms on these small CPU workloads. For example, the SciFact exact
reference takes \SI{2.5774}{\second}, while the selected PLAID points take
\SI{14.5497}{\second} and \SI{16.6456}{\second}.

\subsection{Post-hoc full-score sensitivity}

After observing the shift, a full-corpus score budget was run once per dataset
to determine whether $F=100/400$ alone caused the loss. Because this choice
uses knowledge of held-out behavior, it is an explanatory sensitivity analysis
and not a new selected operating point.

\begin{table}[tb]
  \centering
  \caption{Post-hoc full-score sensitivity. Exact and PLAID medians come from
    the same interleaved run. PLAID/exact is therefore a valid within-run
    online-latency ratio, but the PLAID configuration itself was chosen after
    observing held-out results.}
  \label{tab:plaid-full}
  \small
  \begin{tabular}{@{}l S[table-format=4.0] S[table-format=1.4] S[table-format=2.4] S[table-format=2.2] S[table-format=1.4] S[table-format=1.4]@{}}
    \toprule
    Dataset & {$F$} & {Exact med. (\si{\second})} & {PLAID med. (\si{\second})} & {Ratio ($\times$)} & {Exact R@10} & {Qrels R@10} \\
    \midrule
    SciFact  & 5183 & 3.2697 & 80.6144 & 24.65 & 0.6750 & 0.7850 \\
    NFCorpus & 3633 & 1.2836 & 46.0245 & 35.86 & 0.5975 & 0.1302 \\
    \bottomrule
  \end{tabular}
\end{table}

The larger budget improves exact-ranking recovery but does not reproduce exact
MaxSim, showing that loss also occurs in earlier PLAID stages or compressed
scoring. On this scale the full-score arms remain 24.65--35.86 times slower
than their same-run compiled exact references. This does \emph{not} establish
that PLAID is generally ineffective: PLAID is designed for much larger
collections, and this experiment is limited to one pinned CPU implementation,
four-bit configuration, and two small BEIR corpora. It does establish that no
PLAID speed or quality advantage appears under the controlled conditions used
for this report.

\begin{figure}[tb]
  \centering
  \includegraphics[width=\textwidth]{fig9_controlled_plaid.png}
  \caption{Controlled PLAID results. Selected operating points were frozen on
    the first 10 SciFact queries before held-out evaluation; the full-corpus
    score-budget points are explicitly post hoc. Quality is exact recall@10,
    and latency covers resident normalized query vectors through final ranked
    document IDs. PLAID is slower and less accurate than compiled exact MaxSim
    and the IVF frontiers at this tested CPU scale.}
  \label{fig:plaid}
\end{figure}

Offline PLAID index construction took \SI{1415.73}{\second} on SciFact and
\SI{527.82}{\second} on NFCorpus, producing 182.3 and 131.0 MB indices,
respectively. These costs are excluded from the online latencies above and are
reported separately rather than amortized into a speedup claim.
